\documentclass{article}

\usepackage{../handouts}

\title{CSCI 432 Handout 08: Dynamic Programming}
\author{Name: \rule{3in}{0.15mm} \\~\\
    Collaborators: \rule{3in}{0.15mm} }
\date{17 February 2026}



\begin{document}
\maketitle

\noindent
\section*{The Problem}

Suppose we are given a set $L$ of $n$ line segments in the plane, where each
segment has one endpoint on the line $y=0$ and one endpoint on the line $y=1$,
and all $2n$ endpoints are distinct. Describe and analyze an algorithm to
compute the largest subset of $L$ in which no pair of segments intersects.

From JE, Chapter 3, Problem 19 (a).

\section*{Practice}

\begin{enumerate}
    \item Before getting started, let's consider an example.  What would the desired output be
        for the following input set?
        \begin{center}
            \includegraphics{lines}
        \end{center}
    ~\practice

    \newpage

    \item {\bf Formulate the problem recursively: Specification}. The first step
        in thinking of a dynamic program is to formulate the problem
        recursively. Before even doing that, write down precisely WHAT problem
        you are solving.  Namely, what is the input, the output, and the
        relationship between them?
        \practice

    \item {\bf Formulate the problem recursively: Solution}. Provide a clear
        recursive formula that solves this problem.
        \practice

    \item {\bf Build solution: Identify the subproblems}. What are all of the
        different subproblems that your recursive formula would require solving
        as intermediate steps? Be sure to include the initial input as one of
        those subproblems!
    \practice

    \item {\bf Build solution: Choose a memoization data structure.} What data
        structure can store EVERY subproblem that you identified?
    \practice

    \item {\bf Build solution: identify dependencies}. Draw the dependency
        graph among your subproblems.
    \practice

    \item {\bf Build solution: find a good evaluation order}. Keeping your
        dependency graph in mind, select a good evaluation order so that by the
        time you need a subproblem, it is already computed.
    \practice

    \item {\bf Build solution: analyze space and running time}. What are the
        space and worst-case time complexities of your algorithm?
    \practice

    \item {\bf Write down the algorithm}. Give pseudocode!
    \practice
\end{enumerate}

\end{document}
